% --- Part 5: Problemática y Propuesta ---
\section{Problemática y Propuesta}
\SectionFrame{Problemática y Propuesta}

\begin{frame}{Limitaciones Actuales}
    A pesar de la madurez de CKKS, existe una brecha crítica entre la teoría y la aplicación real en RNAs:
    \vspace{0.3cm}
    \begin{itemize}
        \setlength{\itemsep}{12pt}
        \item \textbf{Limitado por CPU:} Las bibliotecas de CKKS más robustas son extremadamente lentas al escalar a modelos complejos debido a su dependencia de la CPU.
        \item \textbf{Incompatibilidad de GPU:} Herramientas actuales aceleradas por GPU no usan CKKS, limitando la precisión de los modelos.
        \item \textbf{Falta de herramientas:} Ausencia de bibliotecas que permitan construir modelos de redes neuronales abstrayendo la complejidad técnica de FHE para el desarrollador.
    \end{itemize}
\end{frame}

\begin{frame}{Objetivos}
    \textbf{Objetivo General:}
    \begin{itemize}
        \item Desarrollar una biblioteca de código abierto en Python para el cómputo de operaciones tensoriales, utilizando el esquema CKKS y aprovechando la GPU.
    \end{itemize}
    \vspace{0.2cm}
    \textbf{Objetivos Específicos:}
    \begin{itemize}
        \item Diseñar una interfaz en Python para interactuar con FIDESlib y OpenFHE.
        \item Implementar módulos para operaciones de tensores y matrices, abstrayendo la complejidad de CKKS.
        \item Crear casos de uso para inferencia privada en modelos de aprendizaje automático.
        \item Comparar el rendimiento frente a otras bibliotecas de cifrado homomórfico.
    \end{itemize}
\end{frame}

\begin{frame}{Tecnologías}
    \begin{itemize}
        \setlength{\itemsep}{10pt}
        \item \textbf{Python y CUDA:} Lógica de alto nivel en Python y optimización de operaciones no nativas en CUDA.
        \item \textbf{FIDESlib:} Implementación eficiente del esquema CKKS acelerada por GPU.
        \item \textbf{OpenFHE-Python:} Puente principal (\textit{bindings}) para la gestión de las primitivas criptográficas.
        \item \textbf{Pybind11:} Creación de \textit{bindings} personalizados para la comunicación eficiente entre C++ y Python.
    \end{itemize}
\end{frame}

\begin{frame}{Evaluación y Comparativa}
    Se comparará la biblioteca propuesta frente a \textbf{TenSEAL} (CKKS en CPU) y \textbf{Concrete ML} (TFHE en GPU) utilizando modelos de referencia (MLP y CNN) bajo los siguientes criterios:
    \vspace{0.3cm}
    \begin{itemize}
        \setlength{\itemsep}{12pt}
        \item \textbf{Precisión Numérica:} Medir el impacto de las aproximaciones polinomiales y la gestión del ruido frente a modelos en texto plano. Se evaluará la efectividad de los modelos dado el error introducido por las operaciones cifradas.
        \item \textbf{Rendimiento computacional:} Evaluación del tiempo de ejecución y el uso de memoria (RAM/VRAM) para cuantificar la eficiencia de la aceleración por GPU.
    \end{itemize}
    \vfill
    \begin{center}
        \small La biblioteca será publicada de forma abierta en un repositorio público (GitHub).
    \end{center}
\end{frame}
